The diagonal factor can be computed directly.  For
$Q_B=\sum_{i=1}^3 a_iX_i^2$, the one-dimensional Weil-index identity
\[
 \frac{\gamma(ua_i)}{\gamma(a_i)}
       =\frac{\gamma(u)}{\gamma(1)}(u,a_i)_p
\]
and multiplication under orthogonal sums give, relative to the
ordinary theta multiplier cubed, the factor
$(u,a_1a_2a_3)_p$.  These identities are \cite[II.25 and II.28]{weil};
the latter obtains the Hilbert symbol from the four-dimensional
quaternion norm form.  For $B=(-a,-b)$ the norm on $B^0$ is
$aX^2+bY^2+abZ^2$, whose determinant is $(ab)^2$.
There is therefore no additional determinant character for $Q_B$.
\paragraph{The normalized diagonal factor.}
The Weil-index identity used here can be derived without choosing
Gauss-sum signs. The norm form of the quaternion algebra $(u,a)$
is $\langle1,-u,-a,ua\rangle$. By multiplicativity under orthogonal
sums and $\gamma(-q)=\gamma(q)^{-1}$, its index is
\[
 \frac{\gamma(1)\gamma(ua)}{\gamma(u)\gamma(a)}.
\]
For a split quaternion algebra the norm form is hyperbolic and its
index is one; for the division algebra its index is minus one by
\cite[II.28, Proposition 4]{weil}. These are exactly the two values
of $(u,a)_p$. Solving this equality gives
\[
 \frac{\gamma(ua)}{\gamma(a)}
   =\frac{\gamma(u)}{\gamma(1)}(u,a)_p.
\]
Applying it to three diagonal entries gives the ordinary
one-dimensional theta factor cubed times $(u,\det Q_B)_p$.
This derivation includes $p=2$. A change of diagonalization changes
the determinant only by a square, so it changes no symbol.
For $Q_B/m$ the determinant becomes $m^{-3}\det Q_B$, whose
squareclass is $m$. Thus the extra factor is precisely $(u,m)_p$.
This establishes the determinant-character contribution in (O.5)
with the same additive character as the original theta multiplier.

